% TOP_PROBLEM: {"id": 8000004, "title": "Painlev\\'e equations", "category_id": 16, "category": {"id": 16, "name": "physics", "display_name": "Mathematical Physics", "description": "Problems at the intersection of mathematics and physics.", "slug": "physics", "order_index": 16, "created_at": "2026-07-31T15:26:25.671Z"}, "status": "open", "statusQualification": "The source expressly calls this a large-scale project. Its example asks about a Painlevé II solution asymptotic to the Airy function as $x\\to+\\infty$ and the location of complex poles; it is not a claim about all Painlevé solutions simultaneously."}
% FIRST_POSTED: 2026-10-11
% ENTRY_KIND: statement_only
% SUBMISSION: {"kind": "statement_only", "category": "top_problems", "problem_id": "8000004", "proposal_type": "established", "source": "Percy Deift, Some open problems in random matrix theory and the theory of integrable systems (2007), arXiv:0712.0849, Problem 4. https://arxiv.org/abs/0712.0849", "status_evidence": "Inherited source-catalogue status open preserved. Primary source accessed 11 October 2026; compared against the supplied catalogue excerpt. Historical wording and partial-result remarks were checked for scope; no comprehensive current-literature or proof-verification audit was performed.", "status_checked_on": "2026-10-11", "duplicate_check": "Full 15,773-record source registry checked for identical normalized statements; available site registry checked for identical normalized title/statement and canonical-ID collision. Canonical ID is absent from the site registry. Related canonical records 8100006. No mathematical-equivalence literature audit performed.", "attribution": "Percy Deift. Statement preparation: OpenAI GPT-6 Codex assistant, generation 11 October 2026; source-matching and notation editing by the assistant, no independent human mathematical verification. Inherited rights/provenance note: Public arXiv source TeX; arXiv reuse license NEEDS_REVIEW"}
% SOURCE_STATEMENT: "What I have in mind here is not a specific problem, but a project, a very large scale project. The six (nonlinear) Painlev\\'e equations form the core of ``modern special function theory.'' The role that the classical special functions, such as the Airy, Bessel, and Legendre functions, started to play in the 19$^th$ century, has now been greatly expanded by the Painlev\\'e functions. Increasingly, as nonlinear science develops, people are finding that the solutions to an extraordinarily broad array of scientific problems, from neutron scattering theory, to PDEs, to transportation problems, to combinatorics,..., can be expressed in terms of Painlev\\'e functions. What is needed is a project, similar to the Bateman project, or a new volume of Abramowitz and Stegun, devoted to the Painlev\\'e equations. Much can be, and has been, proved regarding the algebraic and asymptotic properties of Painlev\\'e functions. Here the role of integral representations and the classical steepest descent method in deriving precise asymptotics and connection formulae for the classical special functions is played, and expanded, by a Riemann-Hilbert representation of the Painlev\\'e equations, together with the non-commutative steepest descent method introduced in 1993. Very little is known, however, beyond ad hoc calculations, about the numerical solution of the Painlev\\'e equations. If $u(x)$ is the solution of the Painlev\\'e II equation, say, which is asymptotic to the Airy function $Ai(x)$ as $x\\to+\\infty$, one would like to know, for example, the location of its poles in the complex $x$-plane. A modern ``Bateman Project: Painlev\\'e equations'' would not/should not provide tables for such solutions. Rather, it should provide reliable, easy to use software to compute the solutions. Writing useful software for such nonlinear equations presents many challenges, conceptual, philosophical and technical. Without the help of linearity, it is not at all clear how to select a broad enough class of ``representative problems.'' The software should be in the form of a living document where new numerical problems can be addressed by a pool of experts as they arise. And at the technical level, how does one combine asymptotic information about the solutions obtained from the Riemann-Hilbert problem, together with efficient numerical codes in order to compute the solution $u(x)$ at finite values of $x$? I believe that the importance of the Painlev\\'e Project will only grow with time. It should be viewed as creating a national resource and should probably be funded and led at the national level. The NIST Project ``Digital Library of Mathematical Functions'', where Peter Clarkson has a contribution on Painlev\\'e functions, is an encouraging first step."
% Preparation and final self-audit: CONTRIBUTING.md prompts C and D; source/audit artifacts retained in the local physics_statements_2026_10_11 batch.
% Imported from: physics_statements.tex; UnsolvedMath ID 8000004
% !TeX program = lualatex
\documentclass[11pt]{article}
\usepackage[margin=25mm]{geometry}
\usepackage{fontspec}
\setmainfont{Latin Modern Roman}
\usepackage{amsmath,amssymb,mathtools}
\usepackage{unicode-math}
\setmathfont{Latin Modern Math}
\usepackage{xurl,hyperref}
\hypersetup{colorlinks=true,urlcolor=blue}
\setcounter{secnumdepth}{0}
\setlength{\parindent}{0pt}
\setlength{\parskip}{5pt}
\setlength{\emergencystretch}{3em}
\newcommand{\sref}[2]{\hyperlink{source:#1:#2}{[S#2]}}
\newcommand{\cataloguescope}[1]{\paragraph{Recorded target.}#1}
\newcommand{\R}{\mathbb R}
\newcommand{\C}{\mathbb C}
\newcommand{\N}{\mathbb N}
\newcommand{\Z}{\mathbb Z}
\newcommand{\Q}{\mathbb Q}
\newcommand{\F}{\mathbb F}
\newcommand{\eps}{\varepsilon}
\begin{document}
% UnsolvedMath ID 8000004; source catalogue code AMR-079-0004
\setcounter{section}{8000003}
\section{Painlev\'e equations}\label{8000004}
{\small UnsolvedMath ID 8000004 · Mathematical Physics}
\subsection{Definitions and mathematical statement}

\paragraph{Shared notation.}$\mathbb N=\{1,2,\ldots\}$; the standard number fields have their usual meanings. All additional conventions are stated locally.


\paragraph{Definitions and notation.}The Painlevé equations are six families of nonlinear second-order ordinary differential equations distinguished by absence of movable branching singularities. For example, Painlevé II has the form $u''=2u^3+xu+\alpha$. Its solutions are selected by initial or asymptotic data; a pole is a point where the meromorphic solution diverges. Riemann--Hilbert representations encode selected solutions in matrix jump data.

\paragraph{Catalogue wording retained for traceability.}
What I have in mind here is not a specific problem, but a project, a very large scale project. The six (nonlinear) Painlev\'e equations form the core of ``modern special function theory.'' The role that the classical special functions, such as the Airy, Bessel, and Legendre functions, started to play in the 19$^th$ century, has now been greatly expanded by the Painlev\'e functions. Increasingly, as nonlinear science develops, people are finding that the solutions to an extraordinarily broad array of scientific problems, from neutron scattering theory, to PDEs, to transportation problems, to combinatorics,..., can be expressed in terms of Painlev\'e functions. What is needed is a project, similar to the Bateman project, or a new volume of Abramowitz and Stegun, devoted to the Painlev\'e equations. Much can be, and has been, proved regarding the algebraic and asymptotic properties of Painlev\'e functions. Here the role of integral representations and the classical steepest descent method in deriving precise asymptotics and connection formulae for the classical special functions is played, and expanded, by a Riemann-Hilbert representation of the Painlev\'e equations, together with the non-commutative steepest descent method introduced in 1993. Very little is known, however, beyond ad hoc calculations, about the numerical solution of the Painlev\'e equations. If $u(x)$ is the solution of the Painlev\'e II equation, say, which is asymptotic to the Airy function $Ai(x)$ as $x\to+\infty$, one would like to know, for example, the location of its poles in the complex $x$-plane. A modern ``Bateman Project: Painlev\'e equations'' would not/should not provide tables for such solutions. Rather, it should provide reliable, easy to use software to compute the solutions. Writing useful software for such nonlinear equations presents many challenges, conceptual, philosophical and technical. Without the help of linearity, it is not at all clear how to select a broad enough class of ``representative problems.'' The software should be in the form of a living document where new numerical problems can be addressed by a pool of experts as they arise. And at the technical level, how does one combine asymptotic information about the solutions obtained from the Riemann-Hilbert problem, together with efficient numerical codes in order to compute the solution $u(x)$ at finite values of $x$? I believe that the importance of the Painlev\'e Project will only grow with time. It should be viewed as creating a national resource and should probably be funded and led at the national level. The NIST Project ``Digital Library of Mathematical Functions'', where Peter Clarkson has a contribution on Painlev\'e functions, is an encouraging first step.
\paragraph{Scope and source qualification.}The source expressly calls this a large-scale project. Its example asks about a Painlevé II solution asymptotic to the Airy function as $x\to+\infty$ and the location of complex poles; it is not a claim about all Painlevé solutions simultaneously.
\subsection{Short English statement}
Create reliable mathematical and numerical resources for Painlevé functions, including pole locations, asymptotics and connection formulae.
\subsection{Sources}
\begin{itemize}
\item[\hypertarget{source:8000004:1}{S1}] ULAM.AI and the UnsolvedMath Contributors, UnsolvedMath, record 8000004 (AMR-079-0004). Catalogue curation under CC BY 4.0; definitions and English exposition added locally.
\url{https://huggingface.co/datasets/ulamai/UnsolvedMath}.
\item[\hypertarget{source:8000004:2}{S2}] Percy Deift, Some open problems in random matrix theory and the theory of integrable systems (2007), arXiv:0712.0849, Problem 4.
\url{https://arxiv.org/abs/0712.0849}.
\end{itemize}
\label{end:8000004}
\end{document}
